\par\bigskip
\begin{center}
\textbf{VIII. Die Jahre 1813 bis 1827.}\\[.8em]
\textls[100]{a) Die 1813 bis 1816 veröffentlichten Abhandlungen.}
\end{center}
\noindent\hspace*{\parindent}Die Jahre 1813 bis 1816 bringen eine Reihe von Veröffentlichungen, die
mit unserem Gegenstande in mehr oder minder naher Beziehung stehen. Zunächst
ist die bereits genannte \textit{Theoria attractionis}, die dem Jahre 1813 angehört,
zu erwähnen, aus der besonders die im ersten Teile enthaltenen sechs
Theoreme über Integrale für uns von Bedeutung sind. Sie geben die Reduktion
des dreifachen Volumintegrals auf ein Oberflächenintegral und die Auswertung
des die scheinbare Grösse einer geschlossenen Fläche darstellenden
Integrals
\begingroup\makeatletter\tagsleft@true\let\veqno\leqno\makeatother
\begin{equation*}
\int\frac{ds\cos(M,Q)}{r^2},\tag{*}
\end{equation*}
\endgroup
\noindent wo \(ds\) das Flächenelement, \(r\) den Abstand desselben vom Aufpunkte, \(M\) die
Richtung nach diesem hin, \(Q\) die Richtung der nach aussen weisenden Normalen\par
\HistoricalFoot{}{20\textsuperscript{*}}
\clearpage\SourceLeaf{241}\ArithmeticHead{156}{L}{L. SCHLESINGER, ÜBER GAUSS' FUNKTIONENTHEORETISCHE ARBEITEN.}
\noindent bedeutet, und das Integral über die ganze geschlossene Fläche zu erstrecken
ist. Das Integral erweist sich, je nachdem der Aufpunkt ausserhalb,
innerhalb oder auf der Fläche selbst gelegen ist, gleich \(0\), \(-4\pi\) oder \(-2\pi\).
Die Bedeutung dieser Sätze für Potential- und Funktionentheorie ist bekannt,
es ist aber sehr wahrscheinlich\footnote{Vergl. \textsc{Cl. Schaefer}, Werke XI, 2, Abhandlung 2, S. 96—103.}, dass \textsc{Gauss} zu jener Zeit auch schon eine
Reihe von Folgerungen aus diesen Sätzen gezogen hat, die sich erst in der
Abhandlung \textit{Allgemeine Lehrsätze über die im verkehrten Verhältnis des Quadrats
der Entfernung wirkenden Kräfte} von 1839 (Werke V, S. 195) finden, und zwar
nicht nur, wie dort, für das räumliche Potential, sondern auch für das ebene
oder logarithmische, dessen Zusammenhang mit der Theorie der Funktionen
einer komplexen Veränderlichen ihm 1813 sicher längst geläufig war. Finden
wir doch z. B. auf dem Titelblatt des 1802 begonnenen Handbuchs 17, Bc
die Aufzeichnung: »Das Integral von
\[
\frac{\partial^2P}{\partial x^2}+\frac{\partial^2P}{\partial y^2}=0
\]
ist \(P=f(x+yi)+f'(x-yi)\)«\footnote{Über das Auftreten dieses Zusammenhangs bei \textsc{D'Alembert} und \textsc{Euler} vergl. z. B. P. \textsc{Stäckel}, Bibliotheca mathem. (3) 2, 1900, S. 117.}.
\par
Im Jahre 1814 folgt die \textit{Methodus nova integralium valores per approximationem
inveniendi}\footnote{Der G. d. W. vorgelegt am 16. September 1814, Comm. soc. reg. sc. Gott. rec. vol. III, 1816, Werke III, S. 163, \textit{Anzeige} ebenda, S. 202.}, die sich auf die sogenannte mechanische Quadratur von
Integralen der Form \(\displaystyle\int_{-1}^{+1}y\,du\) bezieht, wenn \(y\) für \(-1\leqq u\leqq1\) nach positiven
ganzen Potenzen von \(u\) entwickelbar ist. Während man früher (\textsc{Newton},
\textsc{Cotes}) die Approximation von \(y\) durch eine ganze rationale Funktion \(Y\) stets
in der Weise machte, dass man die Abszissenwerte, für die \(Y\) mit \(y\) übereinstimmen
sollte, äquidistant wählte, zeigte \textsc{Gauss} mit Zuhilfenahme der Kettenbruchentwicklungen
aus den \textit{Disquisitiones} von 1812, dass durch Wahl dieser
Abszissen an den \(n+1\) Stellen, für die die \(n+1\)-te Kugelfunktion
\[
P_{n+1}(u)=\mathrm{const.}\;u^{n+1}F(-\tfrac12n,-\tfrac12(n+1),-(n+\tfrac12),u^{-2})
\]
verschwindet, derselbe Grad der Annäherung erreicht wird, wie bei der Anwendung
der doppelten Anzahl äquidistanter Abszissen\footnote{Vergl. den Bericht über die Abhandlung von \textsc{Gauss} und über die an dieselbe anknüpfenden Arbeiten von \textsc{Jacobi}, \textsc{Christoffel} u. A. bei A. \textsc{Voss}, Encyklopädie, Band II, 1, 1, S. 121ff.}. Eine Anwendung\par
\clearpage\SourceLeaf{242}\ArithmeticHead{157}{R}{VIIIa. DER ZWEITE WURZELEXISTENZBEWEIS.}
\noindent auf die numerische Berechnung des Integrallogarithmus findet sich im letzten
Artikel der Abhandlung.
\par
Die ausführliche Anzeige von J. \textsc{Fr. Pfaffs} \textit{Methodus generalis}\footnote{Juli 1815, Werke III, S. 231.} bezeichnet
\textsc{Gauss} in einem Briefe an \textsc{Bessel}\footnote{\textit{Briefwechsel Gauss-Bessel}, S. 213, vergl. auch an \textsc{Olbers}, \textit{Briefwechsel Olbers-Gauss} I, S. 598.} als »nicht sowohl einen Auszug
daraus, als vielmehr die Quintessenz der \textls[100]{Sache}, in einer durchaus verschiedenen
Darstellung«.
\par
Im Dezember 1815 wird (vergl. oben S. 154) die \textit{Demonstratio nova altera}
des Fundamentalsatzes der Algebra, und bald darauf\footnote{Januar 1816, Werke III, S. 57, \textit{Anzeige} ebenda, S. 107.} die für uns besonders
wichtige \textit{Demonstratio tertia} der G. d. W. vorgelegt. In bezug auf die letztere
schreibt \textsc{Gauss} am 27. Januar 1816 an \textsc{Bessel}\footnote{\textit{Briefwechsel Gauss-Bessel}, S. 230, vergl. auch an \textsc{Olbers}, \textit{Briefwechsel Olbers-Gauss} I, S. 616, und die einleitenden Sätze der Abhandlung selbst.}, dass er diesen dritten Beweis
»erst vor kurzem gefunden« habe, und in einem späteren Briefe an denselben
Adressaten (14. Juni 1816)\footnote{\textit{Briefwechsel Gauss-Bessel}, S. 240.} begleitet er die Zusendung der Separatabdrücke
der beiden Abhandlungen mit folgenden Worten: »In der ersten ist manches,
dem die Leser nicht ansehen werden, wie viel Nachdenken es mich gekostet
hat; dies ist nicht so bei der zweiten, die, nachdem die Hauptidee einmal geboren
war, sogleich in einem Gusse entstand«.
\par
Der Grundgedanke des zweiten Beweises besteht darin, dass man, ausgehend
von einer Gleichung \(m\)-ten Grades mit reellen Koeffizienten
\[
Y=x^m-L_1x^{m-1}+\cdots\pm L_m=0,
\]
eine Gleichung vom Grade \(\tfrac12m(m-1)\) mit ebenfalls reellen Koeffizienten
\(Y_1=0\) bilden kann, deren Wurzeln symmetrische Funktionen von je zweien
der Wurzeln von \(Y=0\) sind, so dass, wenn man eine Wurzel von \(Y_1=0\)
kennt, eine Wurzel von \(Y=0\) durch Auflösung einer quadratischen Gleichung
gefunden werden kann. Wenn dann \(m=2^p\mathbin{.}k\) ist, wo \(k\) eine ungerade
Zahl bedeutet, so ist \(\tfrac12m(m-1)=2^{p-1}\mathbin{.}k'\), wo auch \(k'\) ungerade ist, so
dass man auf diese Weise fortfahrend nach genau \(p\) Schritten zu einer
Gleichung von ungeradem Grade mit reellen Koeffizienten gelangt, für die
die Existenz einer reellen Wurzel gesichert ist. Man kann also von dieser
Wurzel ausgehend, einen aus lauter Quadratwurzeln zusammengesetzten Ausdruck\par
\clearpage\SourceLeaf{243}\ArithmeticHead{158}{L}{L. SCHLESINGER, ÜBER GAUSS' FUNKTIONENTHEORETISCHE ARBEITEN.}
\noindent herstellen, der die ursprüngliche Gleichung \(Y=0\) befriedigt. — Dieser
Gedankengang findet sich in den Abhandlungen von \textsc{Euler} und \textsc{Daviet de
Foncenex}, die \textsc{Gauss} in den Artikeln 7, 8 bezw. 10, 11 der \textit{Dissertation} kritisch
bespricht. Der »Hauptvorwurf« gegen diese Beweisversuche besteht
darin, dass — wie \textsc{Gauss} in der \textit{Anzeige} des zweiten Beweises (Werke III,
S. 105) sagt — »man die Sache so genommen hat, als sei bloss die Form der
Wurzeln zu bestimmen, deren Existenz man voraussetzte, ohne sie zu beweisen«.
Der Zweck des zweiten Beweises war nun — wie \textsc{Gauss} an \textsc{Drobisch} schreibt
(Werke X, 1, S. 107) — »hauptsächlich mit, nachzuweisen, was eigentlich diesen
durchaus illusorischen Beweisversuchen fehlt und wie es ergänzt werden müsse,
und dies konnte nicht ohne Ausführlichkeit geschehen«.
\par
\textsc{Gauss} entwickelt in dem Art. 2 der Abhandlung zunächst das Verfahren
zur Bestimmung des grössten gemeinsamen Teilers zweier ganzer Funktionen,
und dann in den Artikeln 3—5 das zur Darstellung der ganzen symmetrischen
Funktionen von \(m\) Unbestimmten \(a_1,a_2,\ldots,a_m\) durch ihre elementaren symmetrischen
Funktionen \(\lambda_1,\lambda_2,\ldots,\lambda_m\), wobei er (Art. 4) das sogenannte lexikographische
Ordnungsprinzip für die Glieder eines Polynoms in beliebig vielen Veränderlichen
einführt, und (Art. 5) den Nachweis liefert, dass die elementaren symmetrischen
Funktionen von \(m\) Unbestimmten in dem Sinne algebraisch unabhängig
sind, dass eine Identität von der Form
\[
\sum_r C_r\lambda_1^{g_1}\lambda_2^{g_2}\ldots\lambda_m^{g_m}=0
\]
nur möglich ist, wenn alle \(C_r=0\) sind. Es folgt im Art. 6 der Begriff der
Determinante (wir sagen jetzt Diskriminante) \(P\) einer ganzen Funktion \(y\) mit
unbestimmten Koeffizienten, die für die Funktion \(Y\) mit bestimmten Koeffizienten
eine bestimmte Zahl wird, und dann in den Artikeln 7—10 der Beweis
des Satzes, dass das Verschwinden dieser Zahl notwendig und hinreichend
dafür ist, dass \(Y\) mit seiner Derivierten einen gemeinsamen Teiler habe, ohne
dass bei diesem Beweise von der Wurzelexistenz Gebrauch gemacht wird.
— Das ist die »tiefere Untersuchung in der Theorie der Elimination«, die,
wie \textsc{Gauss} am Schluss des Art. 11 der \textit{Dissertation} erklärt, angestellt werden
muss, »um die Methode von \textsc{Foncenex} nach allen Richtungen hin vollständig
zu machen«. Nun bildet er (Art. 12) mit den beiden Unbestimmten \(x,u\) und
den \(m\) Unbestimmten \(a_1,a_2,\ldots,a_m\) das Produkt von \(\tfrac12m(m-1)\) Faktoren\par
\clearpage\SourceLeaf{244}\ArithmeticHead{159}{R}{VIIIa. DER ZWEITE WURZELEXISTENZBEWEIS.}
\[
\zeta=\prod_{i<k}(u-(a_i+a_k)x+a_ia_k),
\]
\noindent das als symmetrische Funktion der \(a_1,\ldots,a_m\) als ganze Funktion
\[
\zeta=f(u,x,\lambda_1,\ldots,\lambda_m)
\]
der \(\lambda_1,\ldots,\lambda_m\) geschrieben werden kann, und zeigt (Artikel 14, 15), dass, wenn
die Diskriminante von \(Y\) nicht verschwindet, die Diskriminante der Funktion
\(\tfrac12m(m-1)\)-ten Grades \(F(u,x)=f(u,x,L_1,\ldots,L_m)\) von \(u\) nicht identisch
in \(x\) verschwinden kann. \textsc{Gauss} stellt nun (Artikel 16—18) eine Identität
auf, die wir in etwas vereinfachter Form\footnote{Vergl. J. \textsc{König}, \textit{Einführung in die allgemeine Theorie der algebraischen Grössen}, 1903, S. 195ff.} wiedergeben wollen. — Bedeutet
\(X\) einen reellen Wert, für den die Diskriminante von \(F(u,X)\) einen von Null
verschiedenen Wert hat, und setzt man
\[
u_{ik}=(a_i+a_k)X-a_ia_k,
\]
also
\[
f(u,X,\lambda_1,\ldots,\lambda_m)=\prod_{i<k}(u-u_{ik}),\qquad(i,k=1,2,\ldots,m)
\]
so befriedigt \(u_{ik}\) eine Gleichung vom Grade \(\tfrac12m(m-1)\), deren Koeffizienten
ganz sind in \(X,\lambda_1,\ldots,\lambda_m\). Ist nun \(s_{ik}\) eine ganze und symmetrische Funktion
von \(a_i\) und \(a_k\) mit in \(\lambda_1,\ldots,\lambda_m\) rationalen Koeffizienten, so folgt aus dem
Ansatz
\begingroup\makeatletter\tagsleft@true\let\veqno\leqno\makeatother
\begin{equation*}
s_{ik}=\sum_r\alpha_ru_{ik}^r\qquad
\begin{gathered}
(r=0,1,2,\ldots,\tfrac12m(m-1)-1)\\
(i,k=1,2,\ldots,m;\ i<k)
\end{gathered}\tag{1}
\end{equation*}
\endgroup
\noindent für die Koeffizienten nach der \textsc{Cramer}schen Regel
\[
\alpha_r=\frac{\Delta_r}{\Delta},
\]
wo \(\Delta\) die Determinante der \(u_{ik}^r\), also \(\Delta^2\) nichts anderes ist, als die Diskriminante
von \(f(u,X,\lambda_1,\ldots,\lambda_m)\), und man erkennt leicht, dass diese \(\alpha_r\) symmetrisch
in den \(a_1,\ldots,a_m\), also rational in den \(\lambda_1,\ldots,\lambda_m\) sind. — Setzt man nun für
\[
\upsilon=z^m-\lambda_1z^{m-1}+\cdots\pm\lambda_m=(z-a_1)(z-a_2)\ldots(z-a_m)
\]
die Zerlegung an
\[
\begin{gathered}
\upsilon=(z^2+s_{m-1}z+s_m)(z^{m-2}+s_1z^{m-3}+\cdots+s_{m-2}),\\
s_{m-1}=-(a_i+a_k),\qquad s_m=a_ia_k,
\end{gathered}
\]
\clearpage\SourceLeaf{245}\ArithmeticHead{160}{L}{L. SCHLESINGER, ÜBER GAUSS' FUNKTIONENTHEORETISCHE ARBEITEN.}
\noindent so sind die \(s_1,\ldots,s_m\) offenbar symmetrisch in \(a_i\) und \(a_k\), also in der Form (1)
darstellbar; es sei
\[
s_\nu=\varphi_\nu(u_{ik})\qquad(\nu=1,2,\ldots,m).
\]
Die entsprechende Darstellung für \(s_m\) und \(s_{m-1}\) ergibt sich sehr einfach; setzt
man nämlich
\[
\left(\frac{\partial\zeta}{\partial u}\right)_{\substack{u=u_{ik}\\x=X}}=u',\quad
\left(\frac{\partial\zeta}{\partial x}\right)_{\substack{u=u_{ik}\\x=X}}=x',
\]
und ist \(u'\ne0\), so folgt aus \(x'=-(a_i+a_k)u'\),
\begingroup\makeatletter\tagsleft@true\let\veqno\leqno\makeatother
\begin{equation*}
s_{m-1}=\frac{x'}{u'},\quad s_m=-\left(u_{ik}+\frac{Xx'}{u'}\right),\tag{1\,a}
\end{equation*}
\endgroup
\noindent und hieraus findet man die Darstellungen \(\varphi_{m-1}(u_{ik}),\ \varphi_m(u_{ik})\), indem man durch
Multiplikation mit den Konjugierten von \(u'\) das \(u_{ik}\) aus dem Nenner entfernt.
\par
Hiernach ist nun die Differenz
\[
G=\upsilon-(z^2+\varphi_{m-1}(u)z+\varphi_m(u))(z^{m-2}+\varphi_1(u)z^{m-3}+\cdots+\varphi_{m-2}(u))
\]
durch \(u-u_{ik}\) und folglich durch \(f(u,X,\lambda_1,\ldots,\lambda_m)\) teilbar, also gleich
\[
H=f(u,X,\lambda_1,\ldots,\lambda_m)g(u,z),
\]
wo \(g(u,z)\) ganz in \(u\) und \(z\) mit in \(X,\lambda_1,\ldots,\lambda_m\) rationalen Koeffizienten ist.
Im wesentlichen ist nun \(G=H\) die von \textsc{Gauss} aufgestellte Identität. — Aus
dieser wird dann im Art. 19 wie folgt weiter geschlossen. Wenn die Gleichung
vom Grade \(\tfrac12m(m-1)\)
\[
f(u,X,L_1,L_2,\ldots,L_m)=F(u,X)=0
\]
die reelle oder komplexe Wurzel \(U\) besitzt, so gibt die Identität \(G=H\) für
\(\lambda_1=L_1,\lambda_2=L_2,\ldots,\lambda_m=L_m,u=U\)
\[
Y=(z^2+\varphi_{m-1}(U)z+\varphi_m(U))(z^{m-1}+\varphi_1(U)z^{m-2}+\cdots+\varphi_{m-2}(U)).
\]
Setzt man in den Formeln (1\,a) für \(u_{ik}\) die Lösung \(U\) und entsprechend für
\(u',x'\) die Werte
\[
\left(\frac{\partial F(u,x)}{\partial u}\right)_{\substack{u=U\\x=X}}=U',\quad
\left(\frac{\partial F(u,x)}{\partial x}\right)_{\substack{u=U\\x=X}}=X',
\]
so wird, da die Diskriminante von \(F(u,X)\) nicht verschwindet, also \(U'\ne0\) ist,
\[
\varphi_{m-1}(U)=\frac{X'}{U'},\quad\varphi_m(U)=-\left(U+\frac{XX'}{U'}\right),
\]
\clearpage\SourceLeaf{246}\ArithmeticHead{161}{R}{VIIIa. DER DRITTE WURZELEXISTENZBEWEIS.}
\noindent d. h. \(Y=0\) wird befriedigt durch die Wurzeln der quadratischen Gleichung
\[
z^2+\frac{X'}{U'}z-\left(U+\frac{XX'}{U'}\right)=0
\]
(Werke III, S. 55). Damit ist die Durchführbarkeit des \textsc{Foncenex}schen Gedankens
gesichert, d. h. die Zurückführung der Wurzelexistenz für eine beliebige
Gleichung mit nicht verschwindender Diskriminante auf die für eine
Gleichung ungeraden Grades geleistet. Auf eine nähere Erörterung der Tatsache,
dass eine Gleichung ungeraden Grades mit reellen Koeffizienten eine reelle
Wurzel besitzt, geht \textsc{Gauss} nicht ein, dagegen kann man wohl sagen, dass diese
\textsc{Gauss}sche Abhandlung nicht nur in den Artikeln 2—5 die Grundlage für
einen exakten Aufbau der formalen Algebra geliefert hat, sondern überhaupt
»das erste und klassische Muster war für die folgerichtige Entwicklung der
allgemeinen Theorie der algebraischen Grössen«\begingroup\renewcommand{\thefootnote}{236a)}\footnote{J. \textsc{König} a. a. O. \textsuperscript{236)}, S. 194, vergl. S. 93. Eine Darstellung des zweiten Beweises findet man bei P. \textsc{Gordan}, Mathem. Annalen 10 (1876), S. 572 und \textit{Vorlesungen über Invariantentheorie} I (1885), § 12. Bei \textsc{Kronecker}, \textit{Festschrift}, 1882, § 13, und ihm folgend bei einer Reihe anderer Algebraiker, beginnend mit J. \textsc{König} a. a. O. bis zu der neuesten Arbeit von K. \textsc{Dörge}, Sitzungsberichte Berlin, 1928, S. 87, geht der Darstellung des eigentlichen zweiten Beweises der Nachweis voraus, dass es stets einen Körper \(K'\) gibt, der den Körper \(K\) der Koeffizienten der Gleichung \(Y=0\) enthält, und in dem diese Gleichung eine Lösung besitzt. Dadurch ist das geleistet, was \textsc{Gauss} bei seinen Vorgängern vermisst (vergl. die oben S. 158 zitierte Bemerkung der \textit{Anzeige}), so dass nur noch zu zeigen bleibt, dass jene Lösung in Form einer komplexen Zahl dargestellt werden kann. Vergl. die besonders einfache Wiedergabe des \textsc{Dörge}schen Beweises bei O. \textsc{Perron}, \textit{Algebra} I, 1932, S. 246, 247. Eine Analyse des Begriffs der reellen Wurzel einer algebraischen Gleichung gibt \textsc{Kronecker}, \textsc{Crelles} Journal 101, 1888, S. 347ff.}\addtocounter{footnote}{-1}\endgroup.
\par
In bezug auf den dritten Beweis \textit{Theorematis de resolubilitate functionum
algebraicarum integrarum in factores reales demonstratio tertia} sagt \textsc{Gauss} in der
\textit{Anzeige} (Werke III, S. 107), dass er, nachdem der zweite Beweis »bereits abgedruckt
war \mbox{.\,.\,.}, bei fortgesetzter Beschäftigung mit demselben Gegenstande
das Glück hatte, dasselbe Ziel auf einem ganz neuen Wege zu erreichen \mbox{.\,.\,.}
Die beiden früheren Beweise unterscheiden sich dadurch, dass der erste (1799)
sehr kurz und einfach, aber zum Teil auf geometrische Betrachtungen gegründet
war, der andere dagegen rein analytisch, aber viel komplizierter \mbox{.\,.\,.}
Dagegen ist nun der gegenwärtige dritte, auf gänzlich verschiedenen Prinzipien
beruhende Beweis ebenfalls rein analytisch, übertrifft aber selbst den ersten
\mbox{.\,.\,.} sehr an Einfachheit und Kürze.«\par
\HistoricalFoot{X 2 Abh. 2.}{21}
